\documentclass[10pt,a4paper]{article}

\usepackage[T1]{fontenc}
\usepackage[utf8]{inputenc}
\usepackage{textcomp}\usepackage{lmodern,amsmath,amssymb,array,tabularx,booktabs,xcolor,enumitem}
\usepackage[margin=20mm]{geometry}
\usepackage[unicode,colorlinks=true,urlcolor=black,linkcolor=black]{hyperref}
\usepackage{mdframed}
\usepackage{xcolor}
\usepackage[normalem]{ulem}
\usepackage{kotex}

\setlength{\parindent}{0pt}
\setlength{\parskip}{5pt}
\setlength{\emergencystretch}{2em}
\setlength{\tabcolsep}{4pt}
\renewcommand{\arraystretch}{1.10}
\newcommand{\hd}[1]{\par\vspace{8pt}\textbf{#1}\par\vspace{3pt}}
\newcommand{\circnum}[1]{\textcircled{\scriptsize #1}}
\newcommand{\note}{\textreferencemark\ }
\newcommand{\red}[1]{\textcolor{red}{#1}}
\newcommand{\under}[2]{\underbrace{#1}_{\text{#2}}}
\newcommand{\paper}[1]{\href{#1}{paper}}
\newcolumntype{Y}{>{\raggedright\arraybackslash}X}

%%%%%%%%%%%%%%%%%%%%%%%%%%%%%%%%%%%%%%%%%%%%%%%%%%%%%%
\begin{document}



%%%%%%%%%%%%%%%%%%%%%%%%%%%%%%%%%%%%%%%%%%%%%%%%%%%%%%
\hd{1. Core model: City selection + vehicle deployment + multiple penalties}

\[
\boxed{\begin{aligned}
\min_{y,q,O,U}\quad
&\under{\sum_c F_c y_c}{Entry costs}
+\under{\sum_c K_c q_c}{Vehicle costs}\\
&+\frac{1}{|S|}\sum_s\sum_c\left[
\under{c_c^+O_c^s}{Remaining capacity}
+\under{c_c^-U_c^s}{Foregone contribution margin}
\right].
\end{aligned}}
\]

%%%%%%%%%%%%%%%%%%%%%%%%%%%%%%%%%%%%%%%%%%%%%%%%%%%%%

In $c^+,c^-$, $+,-$ distinguish excess and shortage, rather than the signs of costs. All demand, throughput, and costs use \textbf{the same evaluation period.} \note I suppose that setting the evaluation period is needed, and the model determines the city/vehicle deployment to apply during that period. 
%%%%%%%%%%%%%%%%%%%%move from pg3 to pg1%%%%%%%%%%%%%
\\

\textbf{2. Constraints }

\[
\boxed{\begin{aligned}
0\le q_c\le\bar q_c y_c&\qquad\forall c,\\
\sum_c q_c\le Q,&\\
\kappa_c q_c-M_c p_c^s=O_c^s-U_c^s
&\qquad\forall c,s,\\
y_c\in\{0,1\},\quad q_c\in\mathbb Z_{\ge0},&\\
O_c^s,U_c^s\ge0.&
\end{aligned}}
\]

{\small
\begin{tabularx}{\linewidth}{@{}>{\raggedright\arraybackslash}p{0.32\linewidth}Y@{}}
\toprule
Constraint & Meaning \\ \midrule
$0\le q_c\le\bar q_c y_c$ & Zero vehicles if the city is not entered; up to the city-specific limit if it is entered \\
$\sum_cq_c\le Q$ & Total vehicle limit \\
$\kappa_c q_c-M_c p_c^s=O_c^s-U_c^s$ & The difference between service capacity and robotaxi demand equals unused capacity minus unmet demand. \\
%%%%%%%%%%%%%%%%%%%%%%%%%%%%%%%%%%%%%%%%%%%%%
%%%%%%%%%%%%%% 삭제 예정%%%%%%%%%%%%%%%%%%%%%%
%%%%%%%%%%%%%%%%%%%%%%%%%%%%%%%%%%%%%%%%%%%%%%

\bottomrule
\end{tabularx}
}


%%%%%%%%%%%%%%%%%%%%%%%%%%%%%%%%%%%%%%%%%%%%%

\hd{3. Explanations}

\textbf{A. Construct two types of demand }

\[
D_c^s = M_c p_c^s,\qquad 
D_c^{N,s}= M_c (1-p_c^s)
\]

{\small
\begin{tabularx}{\linewidth}{@{}>{\raggedright\arraybackslash}p{0.12\linewidth}>{\raggedright\arraybackslash}p{0.42\linewidth}Y@{}}
\toprule
Symbol & Meaning & Explanation \\ \midrule
$M_c$ & Number of \textbf{potential ride-hailing usage opportunities} in city $c$ & Robotaxi demand + human-driven taxi demand + other demand \\
$p_c^s$ & Probability of choosing robotaxi for that usage opportunity & Need to be specified with reference service conditions such as fares and waiting times \\
$D_c^s$ & Robotaxi choice demand & Demand \textbf{before} applying vehicle shortages \\
$D_c^{N,s}$ & Demand not choosing robotaxis & Human-driven taxis, other modes, giving up the trip, etc. \\
\bottomrule
\end{tabularx}
}

\textbf{B. Constraints}\\
1) Once demand $D_c^s$ and service capacity $\kappa_c q_c$ are determined, the following two values are calculated based on the deployment outcome; they are not variables that need to be investigated separately. $(x)^+=\max\{x,0\}$. Both values are measured in trips.

\[
\begin{aligned}
%%%%%%%%%%%%%%%%%%%%%%%%%%%%%%%%%%%%%%
O_c^s&=(\kappa_c q_c-D_c^s)^+&&\text{Remaining service capacity},\\
U_c^s&=(D_c^s-\kappa_c q_c)^+&&\text{Unserved potential demand}.
\end{aligned}
\]\\

%%%%%%%%%%%%%%%%%%%%%%%%%2pg%%%%%%%%%%%%%%%%%%%%%%%%%%%%%
\newpage
\hd{4. Notation used in the model}

{\small
\begin{tabularx}{\linewidth}{@{}>{\raggedright\arraybackslash}p{0.1\linewidth}>{\raggedright\arraybackslash}p{0.37\linewidth}Y@{}}
\toprule
Symbol & Meaning & Type and unit \\ \midrule
$C$ & Set of potential candidate cities & Input set; cities \\
$S$ & Set of demand scenarios & Input set; scenarios \\
$y_c$ & 1 if city $c$ is entered, else 0 & Decision variable; binary (0 or 1) \\
$q_c$ & Number of vehicles deployed in city $c$ & Decision variable; vehicles (integer) \\
$\kappa_c$ & Number of trips one vehicle can serve during the period & Scale parameter; trips/vehicle over the evaluation period.  \\
$Q$ & Total number of vehicles available for deployment & Input parameter; vehicles \\
$\bar q_c$ & Upper bound on vehicle deployment in each city & Input parameter; vehicles \\

%%%%%%%%%%%%%% 비용 항목 %%%%%%%%%%%%%%%%%
$F_c$
& Fixed cost of city entry and basic infrastructure 
& Input parameter; currency over the evaluation period
\\

$K_c$
& Per-vehicle preparation, deployment, and holding cost
over the evaluation period
& Input parameter; currency/vehicle over the evaluation period
\\

$c_c^+$
& Additional loss per unit of unused service capacity
& Input parameter; currency/trip of unused capacity
\\

$c_c^-$
& Foregone contribution margin per unit of unserved
potential demand
& Input parameter; currency/unserved trip
\\

$M_c$ & Number of \textbf{potential ride-hailing usage opportunities} in city $c$ & Input parameter; usage opportunities over the evaluation period. \newline data or assumption \\


%%%%%%%%%%%%%%%%%%%%%%%%%%%%%%%%%%%%
$p_c^s$ & Probability of choosing our company's robotaxi for that usage opportunity & Scenario input; dimensionless probability in $[0,1]$ \\
$D_c^s$ & Robotaxi choice demand & Derived input; trips over the evaluation period. \newline $D_c^s= M_cp_c^s$ \\
$D_c^{N,s}$ & Demand not choosing robotaxis (human-driven taxis, other modes, giving up the trip, etc.) & Derived input; trips over the evaluation period. \newline $D_c^{N,s}= M_c(1-p_c^s)$ \\

$O_c^s$ & Unused service capacity
& Outcome variable; trips over the evaluation period. \newline $O_c^s =(\kappa_c q_c-D_c^s)^+$ \\

$U_c^s$ & Unserved robotaxi demand
& Outcome variable; trips over the evaluation period. \newline $U_c^s=(D_c^s-\kappa_c q_c)^+$ \\


\bottomrule
\end{tabularx}
}

\note Demand-related variables consist of the inputs $p_c^s,M_c$ and outcome variables calculated from scenario demand and deployment in the model. For the baseline model, $M_c$ is fixed.
\note Outcome variables are not quantities to be predicted by collecting data for each one separately. They are outcomes calculated from demand and deployment.

%------------------------------------------
\newpage
\textcolor{blue}{\small\textit{
[Moved from \textbf{baseline(main)} file to \textbf{candidates\_(M,p)} file in \textbf{driving factors/Candidates} folder.]
}}

\begin{mdframed}[linecolor=red!60,linewidth=2pt]
\color{gray!60}

\hd{5. Driving factors}
\hd{A. Driving factors of $M$}
The papers below mainly explain \textbf{observed ride-hailing usage}, so they do not include the `unrealized demand' included in the definition of $M$, making direct estimation impossible. However, I suppose the demand patterns provided by these papers lead to constructing `potential travel opportunities $M$' through modeling.

{\small
\begin{tabularx}{\linewidth}{@{}>{\raggedright\arraybackslash}p{0.16\linewidth}>{\raggedright\arraybackslash}p{0.32\linewidth}>{\raggedright\arraybackslash}p{0.3\linewidth}Y@{}}
\toprule
Variable & Supporting paper & Scope established in the paper & Measurement data \\ \midrule
\textbf{Population and population density} & Elkhouly and Alhadidi, \textit{Ride-sharing Determinants} (\paper{https://arxiv.org/abs/2406.11590}); Zhang et al., \textit{Examining spatial heterogeneity\ldots} (\paper{https://arxiv.org/abs/2209.07980}) & Relationship of population and local density with conventional ride-hailing usage. Density effects may vary by area & \href{https://www.census.gov/programs-surveys/acs}{ACS} \\
\textbf{Employment and economic activity} & Elkhouly and Alhadidi (\paper{https://arxiv.org/abs/2406.11590}) & Uses the economically active population and employment environment to explain usage & \href{https://lehd.ces.census.gov/data/}{LEHD/LODES}, ACS \\
\textbf{Income} & Elkhouly and Alhadidi (\paper{https://arxiv.org/abs/2406.11590}) & Analyzes the relationship between income and usage in Chicago & \href{https://data.census.gov/table/ACSDT5Y2024.B19013}{ACS B19013} \\
\textbf{Share of car-free households} & Elkhouly and Alhadidi (\paper{https://arxiv.org/abs/2406.11590}) & Analyzes the relationship between households without cars and usage & \href{https://data.census.gov/table/ACSDT1Y2022.B08201}{ACS B08201} \\
\textbf{Land-use mix, accessibility, and commercial activity} & Zhang et al., \textit{Examining spatial heterogeneity\ldots} (\paper{https://arxiv.org/abs/2209.07980}) & Nonlinear and spatial differences in the urban environment and ridesourcing usage & \href{https://www.epa.gov/smartgrowth/smart-location-mapping}{EPA SLD} \\
\textbf{Restaurant and leisure-facility density} & Loa, Lee, and Circella, \textit{Just another ride-hailing?} (\paper{https://link.springer.com/article/10.1007/s11116-026-10751-9}) & Relationship with \textbf{robotaxi trip counts in SF}. Does not separate the effects of $M$ and $p$ & \href{https://www.census.gov/programs-surveys/cbp.html}{Census CBP}, OpenStreetMap \\

\textbf{Public transport and walking environment} & Zhang et al. (\paper{https://arxiv.org/abs/2209.07980}); Cats et al., \textit{Beyond the dichotomy\ldots} (\paper{https://doi.org/10.1371/journal.pone.0262496}) & Analyzes the local environment and travel-time competitiveness and substitutability with public transport. Not a uniform increase/decrease effect & EPA SLD, \href{https://gtfs.org/documentation/overview/}{GTFS} \\
\textbf{Airport-related market size} & Zhang et al. (\paper{https://arxiv.org/abs/2209.07980}) & Analyzes airport travel as a separate spatial category. \textbf{Does not estimate the direct effect of airport passenger counts} & \href{https://www.faa.gov/airports/planning_capacity/passenger_allcargo_stats/passenger}{FAA passenger data}. Used as a proxy for market size \\
\textbf{Time, day of week, and season} & Elkhouly and Alhadidi (\paper{https://arxiv.org/abs/2406.11590}) & Analyzes differences in usage over time. Specific variables such as holidays are selected through separate design & Dates and times in trip data \\
\textbf{Weather} & Additional: Liu, Jiang, and Chen, \textit{Quantifying the impact of weather on ride-hailing ridership: Evidence from Haikou, China} (2021, \paper{https://doi.org/10.1016/j.tbs.2021.04.002}) & Analyzes the relationship between weather and ride-hailing usage. Effects in U.S. cities require revalidation & \href{https://www.ncei.noaa.gov/cdo-web/}{NOAA CDO} \\
\bottomrule
\end{tabularx}
}

%The determinants of $M$ are city characteristics (original proposal) + temporal/seasonal context (situation) + market activity (outcome: past Uber/Lyft usage) Past usage is a predictor that can be used separately. However, this does not explain ``why demand occurs''; it can be classified as \textbf{outcome data showing the size (scale) of demand.}

%%%%%%%%%%%%%%%%%%%%%%%%%%%%%%%%%%%%%%%%
\newpage
\hd{B. Driving factors of $p$}

{\small
\begin{tabularx}{\linewidth}{@{}>{\raggedright\arraybackslash}p{0.16\linewidth}>{\raggedright\arraybackslash}p{0.32\linewidth}>{\raggedright\arraybackslash}p{0.3\linewidth}Y@{}}
\toprule
Variable & Supporting paper & Scope of application & Required measurement data \\ \midrule
\textbf{Relative fare} & Additional: Yin and Cherchi, \textit{Willingness to pay for automated taxis\ldots} (2024; online 2022, \paper{https://doi.org/10.1007/s11116-022-10319-3}) & Choice experiment between conventional taxis and automated taxis & Offered fares for both + choice outcomes \\
\textbf{Relative waiting time} & Yin and Cherchi (\paper{https://doi.org/10.1007/s11116-022-10319-3}) & Analyzes choices under different service conditions & Offered waiting times for both + choice outcomes \\
\textbf{In-vehicle travel time} & Yin and Cherchi (\paper{https://doi.org/10.1007/s11116-022-10319-3}) & Analyzes choice trade-offs between fares and time & Offered travel times for both \\
\textbf{Trust, price value, and usage habits} & Additional: Wei et al., \textit{Using the Extended Unified Theory\ldots} (2024, \paper{https://www.nature.com/articles/s41599-024-03271-3}) & Study of robotaxi acceptance intention and behavior. \textbf{Survey scales cannot be interpreted directly as choice probabilities} & Individual attitude and habit surveys \\
\textbf{Age and robotaxi-related knowledge} & Yin and Cherchi (\paper{https://doi.org/10.1007/s11116-022-10319-3}) & Choice differences by individual characteristics. Fixing the sign, such as ``younger people always prefer it, is inappropriate & Link individual characteristics to choice data \\
\textbf{User reviews and in-vehicle support features} & Yin and Cherchi (\paper{https://doi.org/10.1007/s11116-022-10319-3}) & Analyzes customer ratings and features such as destination changes and communication with an operator as choice attributes & Service attributes in the choice experiment \\
\textbf{Experience using robotaxis} & Additional: Yin, Cherchi, and Nettle, \textit{Living in a Ci ty Where Automated Taxis are Operating and Using Them\ldots} (2023, \paper{https://doi.org/10.1177/03611981231155181}) & Analyzes preference differences by whether service operates in the residential area and by individual usage experience & Individual usage experience \\

\textbf{Competitiveness of alternatives such as public transport} & Cats et al., \textit{Beyond the dichotomy\ldots} (\paper{https://doi.org/10.1371/journal.pone.0262496}) & Evidence from \textbf{conventional ride-hailing}. Not a robotaxi-specific preference effect & Availability, fares, and times of alternatives \\
\bottomrule
\end{tabularx}
}


\end{mdframed}

\end{document}


